% Generated draft: requires mathematical and source review.
\documentclass[11pt]{article}
\usepackage[margin=25mm]{geometry}
\usepackage{fontspec}
\setmainfont{Latin Modern Roman}
\usepackage{amsmath,amssymb,mathtools,unicode-math}
\setmathfont{Latin Modern Math}
\usepackage{enumitem,needspace,xurl,hyperref}
\setlength{\parindent}{0pt}
\setlength{\parskip}{5pt}
\setlength{\emergencystretch}{3em}
\newcommand{\N}{\mathbb{N}}
\newcommand{\Z}{\mathbb{Z}}
\newcommand{\Q}{\mathbb{Q}}
\newcommand{\R}{\mathbb{R}}
\newcommand{\C}{\mathbb{C}}
\newcommand{\F}{\mathbb{F}}
\newcommand{\A}{\mathbb{A}}
\newcommand{\PP}{\mathbb{P}}
\newcommand{\E}{\mathbb{E}}
\newcommand{\eps}{\varepsilon}
\newcommand{\dd}{\,\mathrm{d}}
\begin{document}
\section*{UnsolvedMath ID: 20000405}
\textbf{Source problem code:} catalogue entry
\subsection{Definitions and mathematical statement}
\textbf{Original question, attributed to the ``1st Tue speaker''.} Integer Ramanujan sums are
\[c_m(x)=\sum_{\substack{a\bmod m\\(a,m)=1}}\exp(2\pi iax/m).\]
The source asks for function-field analogues and expansions of the von Mangoldt function $\Lambda$ and other arithmetic functions in such sums. The terminal extracted ``2'' is identified by the aid as a page footer, not an exponent or part of the question. No speaker's name is supplied.

In the polynomial interpretation let $A=\F_q[T]$, with sums and factorizations over monic polynomials, $|H|=q^{\deg H}$, $\mu(H)$ the polynomial Möbius function, and $\phi(H)=|(A/H)^\times|$. Choose a nontrivial additive character $\psi$ of $\F_q$ and let $e(f)=\psi(\operatorname{coeff}_{T^{-1}}f)$ at infinity. The polynomial Ramanujan sum is
\[c_H(G)=\sum_{\substack{a\bmod H\\(a,H)=1}}e(aG/H).\]
The standard divisor description is $c_H(G)=\sum_{D\mid(H,G)}|D|\mu(H/D)$. Define $\Lambda_{\deg}(F)=\deg P$ when $F=P^r$ for a monic irreducible $P$ and integer $r\ge1$, and zero otherwise, including $F=1$; put $\Lambda_{\log}=(\log q)\Lambda_{\deg}$. These conventions distinguish degree and logarithmic normalizations.

\textbf{Attached unreviewed account and finite expansion.} The aid attributes polynomial Ramanujan sums to Carlitz, a finite identity written $\sum_Gc_H(G)/|G|=-\Lambda_{\log}(H)/\log q$ to Zheng, and a finite Fourier framework with a von Mangoldt example to Andrade--Hanslope. The range of $G$ and precise hypotheses of Zheng's finite sum are absent from the supplied summary, so that identity is not asserted here as an unrestricted or infinite sum.

The aid's explicit fixed-degree claim is: for every finite field, every integer $n\ge1$, and every monic $F$ of degree $n$,
\[\Lambda_{\deg}(F)=\sum_{\substack{E\ \mathrm{monic,\ squarefree}\\\deg E<n}}a_{E,n}c_E(F),\]
\[a_{E,n}=\frac{\mu(E)}{|E|}\sum_{\substack{K\ \mathrm{monic},\ E\text{-smooth}\\\deg K\le n-\deg E-1}}\frac1{|K|}.\]
Here $E$-smooth means every irreducible factor of $K$ divides $E$, with $K=1$ allowed; for $E=1$ it is the only such polynomial. The report says the coefficient has sign $\mu(E)$ and $|a_{E,n}|\le1/\phi(E)$. The degree restriction and squarefree support make this a finite identity; no convergence convention is needed for that formula.

The aid additionally reports divisor, orthogonality and Dirichlet-series formulas and a varying-modulus Hardy-type expansion whose degree blocks converge absolutely although its individual modulus terms do not. Its actual varying-modulus formula, grouping convention and precise convergence statement are not supplied, so it cannot be interpreted as unconditional absolute convergence over individual moduli. These literature and synthesis claims require independent review.
\subsection{Short English statement}
The source asks for polynomial analogues of Ramanujan sums and von Mangoldt expansions. The aid reports an established function-field framework and a finite expansion on monic polynomials of one degree, with explicit squarefree coefficients. It distinguishes degree from logarithmic normalization and block convergence from convergence over individual moduli.
\subsection{Sources}
\begin{itemize}
\item Local UnsolvedMath registry, permanent ID 20000405, source code catalogue entry; \texttt{aim-algebraic-number-theory-notes.json}, notes[56], tag \texttt{question}.
\item AIM workshop \emph{Arithmetic statistics over finite fields and function fields}, Question 27, ``1st Tue speaker'': \url{https://aimath.org/pastworkshops/arithstatffieldproblems.pdf}. No name or original bibliography is supplied. The aid identifies the trailing 2 as a footer; no PDF was consulted.
\item Attached research attempt 1 names Carlitz, Zheng and Andrade--Hanslope without publication titles, dates or locators. Result \texttt{partial\_result}, problem status at run \texttt{solved}, rigor label \texttt{proved}, novelty confidence low. It is a machine-generated, non-peer-reviewed aid requiring independent verification before citation.
\end{itemize}
\end{document}
